\Needspace{11\baselineskip}
\section{Local H\"older continuity in a big background}
\label{sec:big-minimal-holder}

We now combine the stability estimate above $\rho$ with the regularization of \cref{lem:smooth-background-regularization-rate}, whose curvature error is linear in the regularization parameter. As in the proof of \cref{thm:hermitian-holder-elementary}, a small multiple of the strict barrier absorbs the curvature error of the transforms, and a lower bound for the minimizing radius converts the estimate for the transforms into an estimate for the smoothing on compact subsets of $\Omega$. This gives a local H\"older modulus for solutions with minimal singularities.

\begin{theorem}\label{thm:big-minimal-holder}\index{Holder continuity@H\"older continuity!minimal singularities}
Fix $C_0>0$. Let $u\in\PSH(X,\theta)$ satisfy
\[
 \sup_Xu=0,\qquad V_\theta-C_0\leq u\leq V_\theta,
 \qquad\mathrm{MA}_\theta(u)=f\,\mathrm{d}V,
\]
where $0\leq f\in L^p(\mathrm{d}V)$ and $p>1$.
Then $u$ is locally H\"older continuous on $\Omega$, with every exponent
\[
 0<\beta<\frac{2(p-1)}{np+3(p-1)}.
\]
For fixed bounds on $C_0$ and $\|f\|_p$, the H\"older constants are uniform on
each fixed compact subset of $\Omega$.
\end{theorem}

\begin{proof}
Since $u\geq V_\theta-C_0\geq\rho-C_0$, we may subtract a constant from the strict barrier so that $\rho\leq u-1$; it still satisfies $\theta+\ddc\rho\geq\kappa\omega$ for a fixed
$\kappa>0$. Choose
$0<\gamma<(p-1)/(np+p-1)$ and set $\alpha=2\gamma<1$.
Apply \cref{lem:smooth-background-regularization-rate} with $a=u$, and let $A,K,C$ and $\Phi_{c,\delta}$ be the constants and transforms given there. For small $\delta>0$ set
\[
 c=\delta^\alpha,\qquad e=Ac+C\delta,\qquad s=e/\kappa,
 \qquad v_\delta=\frac{\Phi_{c,\delta}+s\rho}{1+s}.
\]
By the curvature bound of \cref{lem:smooth-background-regularization-rate}, $v_\delta\in\PSH(X,\theta)$. The lower bound for $\Phi_{c,\delta}$ in that lemma implies $\rho\leq\Phi_{c,\delta}$ when $\delta$
is small. Thus
\[
 v_\delta\leq\Phi_{c,\delta}\leq R_\delta u\leq0.
\]
In particular $v_\delta\leq V_\theta$, so \cref{thm:big-reference-one-sided-stability} applies with the reference
$V_\theta$ and the comparison potential $v_\delta$. Together with the $L^1$ smoothing estimate, it gives
\[
 \sup_\Omega(v_\delta-u)_+
 \leq C\|(R_\delta u-u)_+\|_1^\gamma\leq C\delta^\alpha.
\]
Fix a compact subset of $\Omega$ on which $\rho\geq-M$. Since $s\leq C\delta^\alpha$
and $u\leq0$, the identity
\[
 \Phi_{c,\delta}-u=(1+s)(v_\delta-u)+s(u-\rho)
\]
gives $\Phi_{c,\delta}-u\leq C_M\delta^\alpha$ there.

At a radius $t$ where the infimum defining $\Phi_{c,\delta}$ is attained, the monotone lower bound yields
\[
 -c\log(t/\delta)\leq\Phi_{c,\delta}-u+K(\delta^2+\delta)
 \leq C_Mc.
\]
Hence $t\geq a_M\delta$ for a fixed $a_M>0$. As in the proof of
\cref{thm:hermitian-holder-elementary}, monotonicity of the smoothing now gives
\[
 R_{a_M\delta}u-u\leq C_M\delta^\alpha,
 \qquad R_{a_M\delta}u-u\geq-K(a_M\delta)^2
\]
on the chosen compact set. Since every small radius has the form $a_M\delta$, we obtain $|R_tu-u|\leq C_Mt^\alpha$ there for all small $t>0$, with constants uniform on that compact set.
\Needspace{8\baselineskip}
We apply this to a slightly larger compact neighborhood, on which also
$\|\nabla R_tu\|\leq C_M/t$. For nearby points $x,y$ at distance $d>0$, the smooth
function $R_tu$ and a short path in that neighborhood give
\[
 |u(x)-u(y)|\leq C_M(t^\alpha+d/t).
\]
The choice $t=d^{1/(1+\alpha)}$ gives the exponent
$\alpha/(1+\alpha)$, and letting $\gamma$ approach its upper endpoint gives
the asserted range. Pairs of points farther apart are handled by the boundedness of $u$ on the compact neighborhood.
\end{proof}

It suffices that the equation hold on $\Omega$. Suppose that
$\mathrm{MA}_\theta(u)=cf\,\mathrm{d}V$ on $\Omega$, where $c>0$ and
$0\leq f\in L^p(X,\mathrm{d}V)$, and that $u$ satisfies the same bounds as in the theorem.
The lower bound $u\geq V_\theta-C_0\geq\rho-C_0$ shows that $u$ is locally
bounded on $\Omega$, so there $\mathrm{MA}_\theta(u)$ is the local bounded product.
Since $Z$ has zero volume and is not charged by $\mathrm{MA}_\theta(u)$, the equation holds on $X$, and the theorem applies with
$cf$ in place of $f$. The uniform estimates then depend on a bound for
$\|cf\|_p$, which includes the constant $c$.

\subsection{Limits of normalized exponential solutions}

For the normalized exponential family $(v_\lambda,c_\lambda)$ of \cref{lem:big-exponential-uniform-type}, the reference interval $[V_\theta-B,V_\theta]$ is uniform in $0<\lambda\leq1$. Stability therefore controls differences on all of $\Omega$, and the local H\"older estimates strengthen the convergence on compact subsets. We fix the nonzero density $0\leq f\in L^p(\mathrm{d}V)$, $p>1$, and use the constants $B,C,m$ of that lemma.

\Needspace{16\baselineskip}
\begin{proposition}\label{prop:big-exponential-clusters}
For the family $(v_\lambda,c_\lambda)$ of
\cref{lem:big-exponential-uniform-type}, the normalized potentials
are relatively compact for uniform convergence of
differences on $\Omega$. As $\lambda\downarrow0$, every cluster limit
$v$ is normalized, has minimal singularities, and satisfies
\[
 \mathrm{MA}_\theta(v)=cf\,\mathrm{d}V,
\]
where $c>0$ is the unique constant of
\cref{thm:big-minimal-normalized}. Moreover,
\[
 c_\lambda\longrightarrow c,
 \qquad c_\lambda\mathrm{e}^{\lambda v_\lambda}f
          \longrightarrow cf\quad\text{in }L^p(\mathrm{d}V).
\]
Every sequence of parameters in $(0,1]$ has a subsequence for which the
potentials converge in $C^\beta$ on each compact subset of $\Omega$, for
all
\[
 0<\beta<\frac{2(p-1)}{np+3(p-1)}.
\]
\end{proposition}

\begin{proof}
By \cref{lem:big-exponential-uniform-type}, all $v_\lambda$ lie in a common reference interval, and their densities are bounded in $L^p$ by $C\|f\|_p$. We may therefore apply
\cref{cor:big-reference-uniform-convergence} with $\eta=V_\theta$.
Given any sequence $\lambda_j\downarrow0$, we pass to a subsequence with
$v_{\lambda_j}\to v$ uniformly in differences on $\Omega$ and in
$L^1(X)$, and with $c_{\lambda_j}\to c_*\in[m,C]$.
The limit satisfies the same reference bounds and is normalized. Off the pole set,
\[
 0\leq1-\mathrm{e}^{\lambda v_\lambda}
 \leq1-\mathrm{e}^{\lambda(V_\theta-B)}\longrightarrow0
\]
as $\lambda\downarrow0$. Multiplying by $f$ and using dominated convergence, we obtain
$\mathrm{e}^{\lambda v_\lambda}f\to f$ in $L^p$ for the whole family.
The compactness corollary therefore identifies the limit equation as
$\mathrm{MA}_\theta(v)=c_*f\,\mathrm{d}V$, and the uniqueness of the constant gives $c_*=c$.
Since every subsequential limit of $c_\lambda$ as $\lambda\downarrow0$ equals $c$, we conclude that
$c_\lambda\to c$, and the stated convergence of the densities follows.

Finally, \cref{thm:big-minimal-holder} gives uniform local $C^\alpha$
bounds for every $\alpha$ below its endpoint, and the uniform limit
inherits these bounds by pointwise convergence. For
$0<\beta<\alpha$, uniform convergence and the elementary interpolation
inequality on smaller coordinate balls give
\[
 [v_{\lambda_j}-v]_{C^\beta}
 \leq C\|v_{\lambda_j}-v\|_\infty^{1-\beta/\alpha}
          [v_{\lambda_j}-v]_{C^\alpha}^{\beta/\alpha}
 \longrightarrow0,
\]
where the H\"older bounds are taken on slightly larger compact neighborhoods. The
uniformly convergent subsequence given by compactness thus converges in $C^\beta$
on every compact subset of $\Omega$ and for every exponent below the endpoint; no further
extraction depending on the exponent is needed.
\end{proof}

\begin{corollary}\label{cor:big-exponential-symmetries}
Suppose a group $G$ of biholomorphisms of $X$ preserves $\theta$ and the measure
$f\,\mathrm{d}V$. Then every $v_\lambda$ and every cluster limit as
$\lambda\downarrow0$ is $G$-invariant. If the normalized zero-exponent
equation has only one $G$-invariant solution with minimal singularities,
then the whole family $v_\lambda$ converges to it as $\lambda\downarrow0$, uniformly in
differences on $\Omega$ and in each local $C^\beta$ topology of the
proposition.
\end{corollary}
\begin{proof}
For $g\in G$, pullback by $g$ preserves $\theta$-plurisubharmonicity, and
$V_\theta\circ g=V_\theta$ by the definition of the envelope. Pullback also
preserves the positive-exponent equation, so its uniqueness gives
$u_\lambda\circ g=u_\lambda$, and the normalization preserves this equality.
Pullback by a fixed smooth biholomorphism is continuous on $L^1(X)$,
since its Jacobian relative to the fixed smooth volume is bounded.
Hence every cluster limit is invariant almost everywhere, and therefore everywhere for the
canonical representatives. This argument does not require $g$ to
preserve the barrier $\rho$ or the set $\Omega$.

If the invariant solution is unique, every cluster
limit is the same potential. Relative compactness then implies
convergence of the whole family for uniform convergence of differences on $\Omega$:
a sequence staying at a fixed positive distance from that potential would
have a convergent subsequence, a contradiction. The interpolation
argument above upgrades this to the asserted local $C^\beta$ convergence.
\end{proof}

The proposition gives compactness and identifies the equation of every cluster limit, and the symmetry criterion is one way to identify the limit itself. In general, \cref{thm:vanishing-exponential-selection} shows that the whole family converges, and it describes the limit through the greatest subsolution of the zero-exponent equation that satisfies a family of integral normalizations determined by how the constant $c$ depends on the density. This holds even when the zero-exponent equation has several normalized solutions.
